\documentclass[11pt,a4paper]{article}
\usepackage[margin=2.5cm]{geometry}
\usepackage{amsmath,amssymb}
\usepackage{booktabs}
\usepackage{graphicx}
\usepackage{hyperref}
\usepackage{natbib}
\usepackage{siunitx}

\title{FOM1 Architecture Comparison: Methodology}
\author{}
\date{}

\begin{document}
\maketitle

% ============================================================
\section{Objective}
% ============================================================

The goal of this study is to determine which machine learning model best predicts the first figure of merit for cooling (FOM1) of organic molecules. FOM1 is a single number that captures how effective a fluid is as a coolant: a higher FOM1 means a molecule transfers heat more efficiently. Because fluid properties change with temperature, FOM1 is evaluated at two operating temperatures --- \SI{40}{\degreeCelsius} and \SI{100}{\degreeCelsius} --- and separate models are trained for each.

Ten different model architectures are compared. These range from classical tree-based methods to deep neural networks, covering a broad spectrum of modern machine learning approaches. All models are evaluated under identical conditions so that performance differences reflect genuine architectural advantages rather than differences in data handling.

% ============================================================
\section{Dataset}
% ============================================================

The dataset consists of organic molecules for which experimental FOM1 values have been measured at both \SI{40}{\degreeCelsius} and \SI{100}{\degreeCelsius}. Each molecule is represented by its SMILES string --- a standardised text notation that encodes molecular structure. For example, the SMILES string \texttt{CCCCCC} represents hexane (a six-carbon chain).

The two temperature datasets are merged so that only molecules with measurements at \emph{both} temperatures are retained. After removing duplicate entries, the final dataset contains the same set of molecules for both temperature targets, allowing direct comparison of model performance across temperatures.

% ============================================================
\section{Molecular Representations}
% ============================================================

Machine learning models cannot directly process a molecular structure; they require numerical input features. Four different strategies are used to convert SMILES strings into numerical representations, each capturing different aspects of molecular information.

\subsection{RDKit Molecular Descriptors}

Molecular descriptors are numerical values that summarise physicochemical properties of a molecule, such as its molecular weight, number of hydrogen bond donors, topological surface area, and measures of shape and flexibility. The open-source cheminformatics toolkit RDKit computes approximately 210 such descriptors from a SMILES string. After removing descriptors that are constant across all molecules (which carry no predictive information) and replacing any invalid values with zero, approximately 200 descriptors remain per molecule.

These descriptors provide a human-interpretable feature set: each number corresponds to a known chemical property, making it possible to understand \emph{why} a model makes a particular prediction.

\subsection{Morgan Fingerprints}

Morgan fingerprints (also called circular fingerprints or ECFP) encode the local chemical environment around each atom in a molecule. The algorithm works by examining each atom and its neighbours within a specified radius (here, radius = 2, meaning up to two bonds away). Each unique local environment is hashed into one of 2048 binary bits. The result is a 2048-dimensional binary vector where each bit indicates the presence or absence of a particular substructural motif.

Unlike descriptors, fingerprints do not correspond to named physicochemical properties. Instead, they capture structural patterns --- for example, the presence of a hydroxyl group adjacent to an aromatic ring. They are widely used in cheminformatics because they encode fine-grained structural detail that global descriptors may miss.

\subsection{Character-Level SMILES Tokens}

SMILES strings can also be treated as sequences of characters, much like words in a sentence. Each character (e.g., \texttt{C}, \texttt{=}, \texttt{(}, \texttt{O}) is assigned a numerical index, and the resulting sequence of indices is fed directly into sequence-processing models. This approach requires no manual feature engineering --- the model must learn which character patterns are relevant for prediction.

A vocabulary is built from all unique characters in the dataset. Any character not seen during training is mapped to a special ``unknown'' token.

\subsection{ChemBERTa Embeddings}

ChemBERTa is a pre-trained language model designed specifically for chemistry. It was trained on approximately 77 million SMILES strings using a masked language modelling objective (similar to how language models like BERT learn to understand text). During pre-training, the model learns general-purpose representations of molecular structure.

When a SMILES string is passed through ChemBERTa, the model produces a 384-dimensional numerical vector (called the \texttt{[CLS]} embedding) that serves as a compact summary of the molecule. These embeddings can then be used as input features for simpler downstream models.

% ============================================================
\section{Cross-Validation Protocol}
% ============================================================

\subsection{Why Cross-Validation?}

When training a machine learning model, it is essential to evaluate its performance on data it has not seen during training. Simply splitting the data into a single training set and test set can produce misleading results if the split happens to be unrepresentative. Cross-validation addresses this by systematically rotating which data is used for training and which is held out for testing.

\subsection{5-Fold Stratified Cross-Validation}

In 5-fold cross-validation, the dataset is divided into 5 equally sized subsets (called \emph{folds}). The model is trained 5 times: each time, one fold is held out as the test set and the remaining four folds are used for training. Every molecule appears in the test set exactly once, and the results from all 5 test folds are combined to give an overall performance estimate.

The splits are \emph{stratified}, meaning that the distribution of FOM1 values is approximately the same in each fold. This is achieved by first dividing the target values into 5 equal-frequency bins (quintiles) and then ensuring each fold contains a proportional number of molecules from each bin. Without stratification, one fold might by chance contain mostly high-FOM1 molecules, making the evaluation unreliable.

Separate fold splits are generated for the \SI{40}{\degreeCelsius} and \SI{100}{\degreeCelsius} targets because the stratification bins differ between temperatures. However, within a given temperature, all ten architectures use the same fold splits, ensuring a fair comparison. The random seed is fixed at 42 for full reproducibility.

% ============================================================
\section{Architectures}
% ============================================================

This section describes the ten model architectures compared in this study. They fall into four broad categories: tree-based models, feed-forward neural networks, sequence-based neural networks, and transformer-based models. A summary of all architectures is provided in Table~\ref{tab:architectures}.

\begin{table}[htbp]
\centering
\caption{Summary of the ten architectures compared. ``Features'' indicates which molecular representation is used as input.}
\label{tab:architectures}
\begin{tabular}{@{}lll@{}}
\toprule
\textbf{Architecture} & \textbf{Features} & \textbf{Category} \\
\midrule
XGBoost (combined) & Descriptors + fingerprints & Gradient-boosted trees \\
XGBoost (descriptors) & Descriptors only & Gradient-boosted trees \\
XGBoost (fingerprints) & Fingerprints only & Gradient-boosted trees \\
MLP (descriptors) & Descriptors & Feed-forward neural network \\
MLP (fingerprints) & Fingerprints & Feed-forward neural network \\
MLP (combined) & Descriptors + fingerprints & Feed-forward neural network \\
ChemBERTa-2 & SMILES (tokenised) & Transformer (fine-tuned) \\
Ridge (embeddings) & ChemBERTa embeddings & Linear regression \\
BiLSTM & SMILES (character-level) & Recurrent neural network \\
Deep Ensemble MLP & Descriptors & Ensemble of neural networks \\
\bottomrule
\end{tabular}
\end{table}

% ----------------------------------------------------------
\subsection{XGBoost (Descriptors + Fingerprints)}
% ----------------------------------------------------------

\subsubsection*{Background}

XGBoost (eXtreme Gradient Boosting) is a tree-based ensemble method. It builds a prediction by combining many simple decision trees, where each successive tree corrects the errors of the previous ones. A decision tree makes predictions by asking a series of yes/no questions about the input features (e.g., ``is the molecular weight greater than 200?'') and following different branches depending on the answers. Individually, each tree is a weak predictor, but by combining hundreds or thousands of them, the ensemble achieves high accuracy.

XGBoost is widely regarded as one of the best-performing algorithms for tabular data (i.e., data organised in rows and columns of numerical features). It is fast, handles missing values naturally, and includes built-in regularisation to prevent overfitting.

\subsubsection*{Implementation}

The input features are the concatenation of standardised RDKit descriptors and Morgan fingerprints, yielding approximately 2200 features per molecule. Standardisation transforms each feature to have zero mean and unit variance, ensuring that features on different scales contribute equally.

Hyperparameters are tuned automatically using randomised search with 200 random configurations, each evaluated via 5-fold inner cross-validation on the training data. The hyperparameter search space includes:
\begin{itemize}
    \item Number of trees: 100 to 2000
    \item Learning rate: 0.01 to 0.2 (controls how much each tree contributes)
    \item Maximum tree depth: 3 to 10 (controls tree complexity)
    \item Minimum child weight: 1 to 5 (controls when to stop splitting)
    \item Subsample ratio: 0.7 to 1.0 (fraction of data used per tree)
    \item Column subsample ratio: 0.7 to 1.0 (fraction of features used per tree)
\end{itemize}

% ----------------------------------------------------------
\subsection{XGBoost (Descriptors Only)}
% ----------------------------------------------------------

This variant uses the same XGBoost algorithm described above but restricts the input features to the $\sim$200 RDKit molecular descriptors only (no fingerprints). By isolating the descriptors, this model tests whether the physicochemical properties captured by descriptors are sufficient for accurate FOM1 prediction, or whether the additional substructural detail provided by fingerprints is necessary.

All other training details --- hyperparameter search space, randomised search procedure, inner cross-validation --- are identical to the combined XGBoost model above. The only difference is the reduced feature set (approximately 200 features instead of approximately 2200).

% ----------------------------------------------------------
\subsection{XGBoost (Fingerprints Only)}
% ----------------------------------------------------------

This variant uses only the 2048-dimensional Morgan fingerprints as input features, excluding all RDKit descriptors. This tests the complementary hypothesis: whether structural motifs encoded by fingerprints alone contain enough information to predict FOM1, without explicit physicochemical property values.

Comparing the three XGBoost variants (combined, descriptors-only, fingerprints-only) provides insight into the relative importance of these two feature types for FOM1 prediction. If, for example, the descriptors-only model performs nearly as well as the combined model, this would suggest that fingerprints add little additional predictive value beyond what is captured by the physicochemical descriptors.

% ----------------------------------------------------------
\subsection{Multi-Layer Perceptron (MLP) --- Three Variants}
% ----------------------------------------------------------

\subsubsection*{Background}

A Multi-Layer Perceptron (MLP) is the simplest form of neural network. It consists of an input layer, one or more hidden layers, and an output layer. Data flows forward through the network: at each hidden layer, the inputs are multiplied by learned weights, summed, and passed through a non-linear \emph{activation function} (here, ReLU --- Rectified Linear Unit --- which simply sets negative values to zero). This non-linearity allows the network to learn complex, non-linear relationships between molecular features and FOM1.

\emph{Dropout} is applied after each hidden layer as a regularisation technique: during training, a random 10\% of neurons are temporarily deactivated at each step, which prevents the network from becoming overly reliant on any single feature and reduces overfitting.

\subsubsection*{Three Variants}

Three MLP variants are compared, differing only in their input features:

\begin{enumerate}
    \item \textbf{MLP (descriptors):} Uses only the $\sim$200 RDKit descriptors. Architecture: Input $\rightarrow$ 256 $\rightarrow$ 128 $\rightarrow$ 64 $\rightarrow$ 1 (with ReLU and 10\% dropout after each hidden layer).

    \item \textbf{MLP (fingerprints):} Uses only the 2048-dimensional Morgan fingerprints. Same architecture as above.

    \item \textbf{MLP (combined):} Uses the concatenation of descriptors and fingerprints ($\sim$2200 features). A wider first layer (512 neurons instead of 256) accommodates the larger input. Architecture: Input $\rightarrow$ 512 $\rightarrow$ 256 $\rightarrow$ 128 $\rightarrow$ 64 $\rightarrow$ 1.
\end{enumerate}

\subsubsection*{Training Details}

All three MLP variants share the same training procedure:
\begin{itemize}
    \item \textbf{Optimiser:} Adam (a widely used gradient descent variant) with learning rate $10^{-3}$ and weight decay $10^{-5}$ (a mild penalty on large weights to reduce overfitting).
    \item \textbf{Loss function:} Mean Squared Error (MSE) --- the average squared difference between predicted and true FOM1 values.
    \item \textbf{Batch size:} 64 molecules per gradient update.
    \item \textbf{Epochs:} Up to 200 passes through the training data.
    \item \textbf{Early stopping:} Training halts if the validation loss does not improve for 20 consecutive epochs. A 15\% subset of the training fold is held out as an internal validation set for this purpose.
\end{itemize}

All input features are standardised (zero mean, unit variance) before training.

% ----------------------------------------------------------
\subsection{ChemBERTa-2 (Fine-Tuned Transformer)}
% ----------------------------------------------------------

\subsubsection*{Background}

Transformers are a class of neural network originally developed for natural language processing (e.g., machine translation and text generation). They process sequences using an \emph{attention mechanism} that allows the model to weigh the importance of different parts of the input when making predictions. Unlike recurrent models (see BiLSTM below), transformers can process all positions in a sequence simultaneously, making them efficient and effective at capturing long-range dependencies.

ChemBERTa-77M-MTR is a transformer model based on the RoBERTa architecture that has been pre-trained on approximately 77 million SMILES strings. During pre-training, the model learned general-purpose representations of molecular structure by predicting masked (hidden) tokens in SMILES strings. This is analogous to how a language model learns grammar and semantics by reading vast amounts of text.

\subsubsection*{Fine-Tuning Strategy}

Fine-tuning adapts the pre-trained model to our specific task (FOM1 prediction) by continuing to train it on our dataset. A small regression head is added on top of the pre-trained backbone:

\begin{center}
\texttt{[CLS] token} $\rightarrow$ Linear(384, 128) $\rightarrow$ ReLU $\rightarrow$ Dropout(0.1) $\rightarrow$ Linear(128, 1)
\end{center}

The \texttt{[CLS]} token is a special token whose output representation summarises the entire input sequence, serving as the molecular embedding.

Training proceeds in two phases:
\begin{enumerate}
    \item \textit{Phase 1 (epochs 0--4):} The pre-trained backbone is \emph{frozen} (its weights are not updated). Only the regression head is trained with a learning rate of $2 \times 10^{-3}$. This allows the head to adapt to the backbone's representations without disrupting them.
    \item \textit{Phase 2 (epochs 5 onward):} The backbone is \emph{unfrozen} and fine-tuned with a much smaller learning rate ($10^{-4}$) to gently adjust its representations, while the head continues training at $2 \times 10^{-3}$. This differential learning rate prevents the pre-trained knowledge from being overwritten.
\end{enumerate}

Training runs for up to 50 epochs with early stopping (patience = 10). Gradient norms are clipped at 1.0 to prevent unstable updates. SMILES strings are tokenised with a maximum length of 128 tokens.

% ----------------------------------------------------------
\subsection{Ridge Regression on ChemBERTa Embeddings}
% ----------------------------------------------------------

\subsubsection*{Background}

Ridge regression is a form of linear regression with an added penalty (called $L_2$ regularisation) that discourages overly large coefficients. The regularisation strength is controlled by a parameter $\alpha$: larger $\alpha$ means more regularisation (simpler model, less prone to overfitting) at the cost of potentially higher bias. Despite its simplicity, ridge regression can be surprisingly effective when the input features are already informative representations.

\subsubsection*{Implementation}

The frozen (non-fine-tuned) ChemBERTa backbone is used to extract 384-dimensional \texttt{[CLS]} embeddings for each molecule. These embeddings capture structural information learned during pre-training. The embeddings are standardised (zero mean, unit variance) and fed into a Ridge regression model.

The regularisation strength $\alpha$ is selected automatically from 50 logarithmically spaced candidates between $10^{-3}$ and $10^{3}$ using built-in cross-validation (\texttt{RidgeCV}).

This model serves as a fast baseline to assess how much of the predictive power of ChemBERTa comes from its pre-trained representations alone (without fine-tuning) and how much benefit fine-tuning provides.

% ----------------------------------------------------------
\subsection{BiLSTM (Bidirectional Long Short-Term Memory)}
% ----------------------------------------------------------

\subsubsection*{Background}

Recurrent Neural Networks (RNNs) are designed to process sequential data by maintaining a hidden state that is updated as each element of the sequence is read. However, standard RNNs struggle with long sequences because gradients (the signals used to update weights) can vanish or explode over many timesteps. The Long Short-Term Memory (LSTM) architecture addresses this with a gating mechanism that controls which information is retained, forgotten, or output at each step, enabling it to capture long-range dependencies.

A \emph{bidirectional} LSTM (BiLSTM) processes the sequence in both directions --- left-to-right and right-to-left --- and concatenates the resulting hidden states. This allows the model to use context from both sides of any position in the sequence, which is important because the significance of a character in a SMILES string often depends on characters that appear both before and after it.

\subsubsection*{Architecture}

\begin{enumerate}
    \item \textbf{Embedding layer:} Each character in the SMILES string is mapped to a 64-dimensional learned vector. This transforms discrete character indices into continuous representations that the network can process.
    \item \textbf{BiLSTM:} Two stacked BiLSTM layers with hidden dimension 128 (256 total when both directions are concatenated). Dropout of 0.1 is applied between layers.
    \item \textbf{Masked mean pooling:} The BiLSTM outputs a vector at every character position. Since SMILES strings vary in length and shorter strings are padded with zeros, only the outputs at valid (non-padded) positions are averaged to produce a single fixed-length molecular representation.
    \item \textbf{Regression head:} Linear(256, 256) $\rightarrow$ ReLU $\rightarrow$ Dropout(0.1) $\rightarrow$ Linear(256, 1).
\end{enumerate}

Training uses Adam (lr = $10^{-3}$, weight decay = $10^{-5}$), MSE loss, gradient clipping at 1.0, batch size 64, up to 100 epochs with early stopping (patience = 15).

% ----------------------------------------------------------
\subsection{Deep Ensemble MLP}
% ----------------------------------------------------------

\subsubsection*{Background}

A deep ensemble trains multiple neural networks independently (with different random initialisations) on the same data and combines their predictions. Because each network converges to a different solution, the disagreement between their predictions provides a natural measure of \emph{uncertainty} --- if all networks agree, the prediction is likely reliable; if they disagree, the model is uncertain.

This is particularly valuable for downstream applications such as Bayesian optimisation, where uncertainty estimates guide the search for new candidate molecules.

\subsubsection*{Architecture and Training}

Five independent MLP members are trained, each with hidden layers of 256, 128, and 64 neurons. Unlike the standard MLP, each member has two output heads:
\begin{itemize}
    \item A \textbf{mean head} $\mu$ that predicts the expected FOM1 value.
    \item A \textbf{log-variance head} $\log \sigma^2$ that predicts the model's own estimate of its prediction uncertainty (\emph{aleatoric} uncertainty --- inherent noise in the data).
\end{itemize}

Each member is trained with the Gaussian negative log-likelihood (NLL) loss:
\begin{equation}
    \mathcal{L} = \frac{1}{2} \left\langle \log \sigma^2 + \frac{(y - \mu)^2}{\sigma^2} \right\rangle
\end{equation}
This loss encourages the model to predict both an accurate mean \emph{and} a well-calibrated variance. If the model predicts a large variance for a given molecule, it is penalised less for errors on that molecule, but penalised for being overconfident (small $\sigma^2$) when the error is actually large.

\subsubsection*{Ensemble Prediction and Uncertainty Decomposition}

The final prediction is an inverse-variance weighted mean across the five members: members that are more confident (smaller $\sigma^2$) contribute more to the final prediction. The total uncertainty is decomposed into two components:

\begin{itemize}
    \item \textbf{Epistemic uncertainty} (model uncertainty) --- the variance of the member means $\mu_m$:
    \begin{equation}
        \sigma^2_{\text{epistemic}} = \mathrm{Var}_{m}[\mu_m] = \frac{1}{M}\sum_{m=1}^{M}(\mu_m - \bar{\mu})^2
    \end{equation}
    This reflects disagreement between ensemble members and is high when the model has insufficient training data in a region of chemical space. It can be reduced by providing more training data.

    \item \textbf{Aleatoric uncertainty} (data uncertainty) --- the average of the member variances:
    \begin{equation}
        \sigma^2_{\text{aleatoric}} = \frac{1}{M} \sum_{m=1}^{M} \sigma^2_m
    \end{equation}
    This reflects inherent noise in the measurements and cannot be reduced by collecting more data.
\end{itemize}

The total uncertainty is:
\begin{equation}
    \sigma_{\text{total}} = \sqrt{\sigma^2_{\text{epistemic}} + \sigma^2_{\text{aleatoric}}}
\end{equation}

Calibration is assessed by checking whether the empirical coverage matches theoretical expectations: approximately 68.3\% of true values should fall within $\pm 1\sigma_{\text{total}}$ of the prediction, and 95.4\% within $\pm 2\sigma_{\text{total}}$.

% ============================================================
\section{Evaluation Metrics}
% ============================================================

All architectures are evaluated on held-out test data using three complementary metrics:

\begin{itemize}
    \item \textbf{Coefficient of determination ($R^2$):} Measures how much of the variance in the true FOM1 values is explained by the model's predictions. An $R^2$ of 1.0 indicates perfect prediction; 0.0 indicates the model is no better than predicting the mean value for every molecule. This is the primary metric used to rank architectures.

    \item \textbf{Root Mean Squared Error (RMSE):} The square root of the average squared prediction error. It is in the same units as FOM1 and gives higher weight to large errors. Mathematically:
    \begin{equation}
        \text{RMSE} = \sqrt{\frac{1}{N}\sum_{i=1}^{N}(y_i - \hat{y}_i)^2}
    \end{equation}

    \item \textbf{Mean Absolute Error (MAE):} The average absolute prediction error. Unlike RMSE, it weights all errors equally and is more robust to outliers:
    \begin{equation}
        \text{MAE} = \frac{1}{N}\sum_{i=1}^{N}|y_i - \hat{y}_i|
    \end{equation}
\end{itemize}

Results are reported in two ways:
\begin{enumerate}
    \item \textbf{Pooled metrics:} All predictions from the 5 test folds are concatenated and metrics are computed on the combined set. This gives the most statistically stable estimate of overall performance.
    \item \textbf{Per-fold metrics:} Metrics are computed separately for each fold, and the mean $\pm$ standard deviation across folds is reported. This captures the variability of performance across different data splits.
\end{enumerate}

% ============================================================
\section{Experimental Design}
% ============================================================

Each of the ten architectures is trained independently at each of the two temperatures (\SI{40}{\degreeCelsius} and \SI{100}{\degreeCelsius}), yielding $10 \times 2 = 20$ model configurations. Each configuration is evaluated across 5 cross-validation folds, resulting in a total of $20 \times 5 = 100$ individual training runs.

All 100 training runs are independent of one another and are executed in parallel as a PBS array job on a high-performance computing (HPC) cluster. Each job is allocated 16 CPU cores and 64\,GB of RAM with a maximum walltime of 72 hours. All neural network models are trained on CPU (no GPU acceleration) to ensure compatibility with the HPC environment.

After all jobs complete, a comparison script aggregates the results and produces:
\begin{itemize}
    \item Per-temperature ranked tables and bar charts of $R^2$ across architectures.
    \item Parity plots (predicted vs.\ true FOM1) for each architecture.
    \item Residual distribution histograms.
    \item Training loss curves for neural network architectures.
    \item A combined cross-temperature comparison to assess whether architecture rankings are consistent between \SI{40}{\degreeCelsius} and \SI{100}{\degreeCelsius}.
\end{itemize}

% ============================================================
\section{Software}
% ============================================================

\begin{itemize}
    \item \textbf{Molecular descriptors and fingerprints:} RDKit (open-source cheminformatics toolkit)
    \item \textbf{Gradient boosting:} XGBoost
    \item \textbf{Neural networks:} PyTorch (CPU mode)
    \item \textbf{Transformer backbone:} HuggingFace Transformers library (ChemBERTa-77M-MTR, DeepChem)
    \item \textbf{Classical machine learning:} scikit-learn (Ridge regression, data standardisation, stratified cross-validation, randomised hyperparameter search)
\end{itemize}

All experiments use Python 3.10 with fixed random seeds for reproducibility.

% ============================================================
\section{Results}
% ============================================================

\subsection{Overall Performance at \SI{40}{\degreeCelsius}}

Table~\ref{tab:results_40} presents the pooled test-set performance of all ten architectures for FOM1 prediction at \SI{40}{\degreeCelsius}, ranked by $R^2$.

\begin{table}[htbp]
\centering
\caption{Architecture comparison for FOM1 at \SI{40}{\degreeCelsius}. Architectures are ranked by pooled $R^2$. Per-fold values report mean $\pm$ standard deviation across 5 folds.}
\label{tab:results_40}
\begin{tabular}{@{}l S[table-format=1.3] S[table-format=2.2] S[table-format=2.2] c@{}}
\toprule
\textbf{Architecture} & {$R^2$} & {RMSE} & {MAE} & {$R^2$ (fold mean $\pm$ std)} \\
\midrule
XGBoost (combined)      & 0.880 &  4.86 &  3.57 & $0.878 \pm 0.023$ \\
XGBoost (descriptors)   & 0.872 &  5.03 &  3.68 & $0.870 \pm 0.026$ \\
MLP (descriptors)       & 0.872 &  5.03 &  3.69 & $0.868 \pm 0.037$ \\
Ridge (embeddings)      & 0.858 &  5.29 &  3.88 & $0.855 \pm 0.041$ \\
ChemBERTa-2             & 0.853 &  5.38 &  4.05 & $0.849 \pm 0.041$ \\
Deep Ensemble MLP       & 0.827 &  5.84 &  4.55 & $0.825 \pm 0.021$ \\
XGBoost (fingerprints)  & 0.526 &  9.67 &  7.41 & $0.522 \pm 0.058$ \\
BiLSTM                  & 0.483 & 10.09 &  7.68 & $0.465 \pm 0.334$ \\
MLP (desc+fp)           & 0.409 & 10.79 &  7.34 & $0.408 \pm 0.067$ \\
MLP (fingerprints)      & {$-0.152$} & 15.07 & 10.97 & $-0.152 \pm 0.112$ \\
\bottomrule
\end{tabular}
\end{table}

The top six architectures all achieve $R^2 > 0.82$, with XGBoost (combined) leading at $R^2 = 0.880$ (RMSE = 4.86). XGBoost (descriptors) and MLP (descriptors) follow closely, both near $R^2 = 0.87$. Ridge regression on ChemBERTa embeddings and fine-tuned ChemBERTa-2 also perform well ($R^2 = 0.858$ and $0.853$ respectively).

A clear performance gap separates the top six from the remaining four architectures. XGBoost (fingerprints) achieves only $R^2 = 0.526$, while BiLSTM shows high variability across folds (std = 0.334). MLP (combined) and MLP (fingerprints) perform poorly, with MLP (fingerprints) yielding a negative $R^2$, indicating predictions worse than a constant-mean baseline.

\subsection{Overall Performance at \SI{100}{\degreeCelsius}}

Table~\ref{tab:results_100} presents the corresponding results at \SI{100}{\degreeCelsius}.

\begin{table}[htbp]
\centering
\caption{Architecture comparison for FOM1 at \SI{100}{\degreeCelsius}. Architectures are ranked by pooled $R^2$.}
\label{tab:results_100}
\begin{tabular}{@{}l S[table-format=1.3] S[table-format=2.2] S[table-format=2.2] c@{}}
\toprule
\textbf{Architecture} & {$R^2$} & {RMSE} & {MAE} & {$R^2$ (fold mean $\pm$ std)} \\
\midrule
Ridge (embeddings)      & 0.823 &  6.07 &  4.36 & $0.822 \pm 0.043$ \\
XGBoost (descriptors)   & 0.814 &  6.23 &  4.39 & $0.812 \pm 0.023$ \\
MLP (descriptors)       & 0.808 &  6.33 &  4.62 & $0.806 \pm 0.036$ \\
XGBoost (combined)      & 0.807 &  6.34 &  4.43 & $0.806 \pm 0.015$ \\
ChemBERTa-2             & 0.766 &  6.99 &  5.07 & $0.764 \pm 0.045$ \\
Deep Ensemble MLP       & 0.679 &  8.19 &  6.42 & $0.675 \pm 0.057$ \\
XGBoost (fingerprints)  & 0.545 &  9.74 &  7.53 & $0.542 \pm 0.045$ \\
BiLSTM                  & 0.312 & 11.98 &  9.53 & $0.289 \pm 0.280$ \\
MLP (desc+fp)           & {$-0.159$} & 15.54 & 11.13 & $-0.187 \pm 0.219$ \\
MLP (fingerprints)      & {$-1.077$} & 20.81 & 14.77 & $-1.143 \pm 0.523$ \\
\bottomrule
\end{tabular}
\end{table}

At \SI{100}{\degreeCelsius}, the ranking shifts notably. Ridge regression on ChemBERTa embeddings takes the top position ($R^2 = 0.823$), overtaking XGBoost (combined), which drops to fourth place ($R^2 = 0.807$). XGBoost (descriptors) and MLP (descriptors) remain strong second and third respectively. All models perform worse at \SI{100}{\degreeCelsius} than at \SI{40}{\degreeCelsius}, with the best $R^2$ dropping from 0.880 to 0.823. The same four architectures that performed poorly at \SI{40}{\degreeCelsius} also fail at \SI{100}{\degreeCelsius}, with MLP (fingerprints) reaching $R^2 = -1.077$.

\subsection{Cross-Temperature Comparison}

Figure~\ref{fig:cross_temp} shows the pooled $R^2$ for all architectures side by side at both temperatures. The top-performing group (XGBoost combined, XGBoost descriptors, MLP descriptors, Ridge embeddings, ChemBERTa-2) is consistent across temperatures, although the internal ranking changes. The performance degradation from \SI{40}{\degreeCelsius} to \SI{100}{\degreeCelsius} is most pronounced for the Deep Ensemble MLP (from 0.827 to 0.679) and XGBoost combined (from 0.880 to 0.807), while Ridge embeddings shows the smallest drop among the top models (from 0.858 to 0.823).

\begin{figure}[htbp]
    \centering
    \includegraphics[width=\textwidth]{figures/combined/cross_temperature_comparison.pdf}
    \caption{Cross-temperature comparison of all ten architectures. Blue bars show pooled $R^2$ at \SI{40}{\degreeCelsius}; red bars at \SI{100}{\degreeCelsius}. Error bars indicate the standard deviation across 5 folds.}
    \label{fig:cross_temp}
\end{figure}

\subsection{Parity Plots}

Figure~\ref{fig:parity_40} shows predicted versus true FOM1 values for all architectures at \SI{40}{\degreeCelsius}. The top-performing models (XGBoost combined, XGBoost descriptors, MLP descriptors, Ridge embeddings, ChemBERTa-2) show tight clustering around the diagonal, with scatter increasing for molecules at extreme FOM1 values. The poorly performing models (MLP fingerprints, MLP combined, BiLSTM) show wide scatter with systematic bias, particularly MLP (fingerprints), which collapses to a narrow band of predictions regardless of the true value.

\begin{figure}[htbp]
    \centering
    \includegraphics[width=\textwidth]{figures/fom1_40/parity_plots.pdf}
    \caption{Parity plots for all architectures at \SI{40}{\degreeCelsius}. Each panel shows predicted versus true FOM1 values pooled across 5 folds. The red dashed line indicates perfect prediction.}
    \label{fig:parity_40}
\end{figure}

\begin{figure}[htbp]
    \centering
    \includegraphics[width=\textwidth]{figures/fom1_100/parity_plots.pdf}
    \caption{Parity plots for all architectures at \SI{100}{\degreeCelsius}.}
    \label{fig:parity_100}
\end{figure}

\subsection{Residual Distributions}

The residual distributions (Figures~\ref{fig:resid_40} and~\ref{fig:resid_100}) show that the top-performing models produce residuals that are approximately centred on zero with relatively narrow spread. The poorly performing models exhibit wide, often asymmetric residual distributions with heavy tails extending beyond $\pm 40$ FOM1 units, confirming that these models fail to capture the underlying structure--property relationship.

\begin{figure}[htbp]
    \centering
    \includegraphics[width=\textwidth]{figures/fom1_40/residual_distributions.pdf}
    \caption{Residual distributions (true $-$ predicted) for all architectures at \SI{40}{\degreeCelsius}.}
    \label{fig:resid_40}
\end{figure}

\begin{figure}[htbp]
    \centering
    \includegraphics[width=\textwidth]{figures/fom1_100/residual_distributions.pdf}
    \caption{Residual distributions for all architectures at \SI{100}{\degreeCelsius}.}
    \label{fig:resid_100}
\end{figure}

\subsection{Training Curves}

The training loss curves for the neural network architectures (Figures~\ref{fig:curves_40} and~\ref{fig:curves_100}) reveal several patterns. MLP (descriptors) converges smoothly within approximately 25 epochs at both temperatures, with training and validation losses closely tracking each other, indicating good generalisation. ChemBERTa-2 shows an initial plateau during the frozen-backbone phase (epochs 0--4), followed by a sharp drop when the backbone is unfrozen, and converges within 30--40 epochs. BiLSTM exhibits noisy training dynamics with substantial variation across folds, consistent with its high fold-to-fold $R^2$ variability. MLP (fingerprints) and MLP (combined) show a large gap between training and validation loss, with validation loss plateauing at a high level, indicating overfitting to the high-dimensional fingerprint input.

\begin{figure}[htbp]
    \centering
    \includegraphics[width=\textwidth]{figures/fom1_40/training_curves.pdf}
    \caption{Training and validation loss curves for neural network architectures at \SI{40}{\degreeCelsius}. Each panel shows all 5 folds (blue: training loss; red: validation loss).}
    \label{fig:curves_40}
\end{figure}

\begin{figure}[htbp]
    \centering
    \includegraphics[width=\textwidth]{figures/fom1_100/training_curves.pdf}
    \caption{Training and validation loss curves for neural network architectures at \SI{100}{\degreeCelsius}.}
    \label{fig:curves_100}
\end{figure}

\subsection{Feature Importance (SHAP Analysis)}
\label{sec:shap_results}

To understand which molecular descriptors drive the XGBoost predictions, we applied SHAP (SHapley Additive exPlanations) analysis using the TreeExplainer algorithm, which provides exact Shapley values for tree-based models. For each of the five cross-validation folds, we computed SHAP values on the held-out test set using the corresponding XGBoost (descriptors) model. Since each molecule appears in exactly one test fold, concatenating across all folds yields SHAP values for the entire dataset without data leakage.

Figures~\ref{fig:shap_beeswarm_40} and~\ref{fig:shap_beeswarm_100} show beeswarm plots of the top 20 descriptors ranked by mean absolute SHAP value at \SI{40}{\degreeCelsius} and \SI{100}{\degreeCelsius}, respectively. Each point represents one molecule, coloured by the feature value (red = high, blue = low), with horizontal position indicating the SHAP contribution to the predicted FOM1 value. The corresponding bar plots (Figures~\ref{fig:shap_bar_40} and~\ref{fig:shap_bar_100}) summarise feature importance as the mean absolute SHAP value across all molecules.

Figure~\ref{fig:shap_cross_temp} provides a direct comparison of the top 15 descriptors by mean $|\text{SHAP}|$ at both temperatures. This reveals which descriptors maintain consistent importance across operating conditions and which become more or less influential as temperature increases.

\begin{figure}[htbp]
    \centering
    \includegraphics[width=\textwidth]{figures/fom1_40/shap_beeswarm.pdf}
    \caption{SHAP beeswarm plot for the XGBoost (descriptors) model at \SI{40}{\degreeCelsius}. Each point represents one molecule; colour indicates the feature value (red = high, blue = low). The top 20 descriptors by mean $|\text{SHAP}|$ are shown.}
    \label{fig:shap_beeswarm_40}
\end{figure}

\begin{figure}[htbp]
    \centering
    \includegraphics[width=\textwidth]{figures/fom1_100/shap_beeswarm.pdf}
    \caption{SHAP beeswarm plot for the XGBoost (descriptors) model at \SI{100}{\degreeCelsius}.}
    \label{fig:shap_beeswarm_100}
\end{figure}

\begin{figure}[htbp]
    \centering
    \includegraphics[width=\textwidth]{figures/fom1_40/shap_bar.pdf}
    \caption{Mean $|\text{SHAP}|$ feature importance for the XGBoost (descriptors) model at \SI{40}{\degreeCelsius}.}
    \label{fig:shap_bar_40}
\end{figure}

\begin{figure}[htbp]
    \centering
    \includegraphics[width=\textwidth]{figures/fom1_100/shap_bar.pdf}
    \caption{Mean $|\text{SHAP}|$ feature importance for the XGBoost (descriptors) model at \SI{100}{\degreeCelsius}.}
    \label{fig:shap_bar_100}
\end{figure}

\begin{figure}[htbp]
    \centering
    \includegraphics[width=\textwidth]{figures/combined/cross_temperature_comparison.pdf}
    \caption{Cross-temperature comparison of descriptor importance. The top 15 descriptors by average mean $|\text{SHAP}|$ are shown, with bars comparing the \SI{40}{\degreeCelsius} and \SI{100}{\degreeCelsius} models.}
    \label{fig:shap_cross_temp}
\end{figure}

% ============================================================
\section{Discussion}
% ============================================================

\subsection{Descriptors Dominate Over Fingerprints}

The most striking result is the dramatic difference between descriptor-based and fingerprint-based models. At \SI{40}{\degreeCelsius}, XGBoost (descriptors) achieves $R^2 = 0.872$ while XGBoost (fingerprints) manages only $R^2 = 0.526$ --- a gap of 0.35. The same pattern holds for the MLP variants: MLP (descriptors) reaches $R^2 = 0.872$ while MLP (fingerprints) yields a negative $R^2$. This pattern is consistent across both temperatures.

This indicates that FOM1 is primarily governed by global physicochemical properties (molecular weight, density-related descriptors, viscosity proxies, thermal conductivity correlates) rather than specific substructural motifs. This is physically intuitive: the figure of merit for cooling depends on bulk thermophysical properties such as density, specific heat capacity, thermal conductivity, and viscosity, all of which are better captured by whole-molecule descriptors than by local structural fingerprints.

The combined XGBoost model (descriptors + fingerprints) only marginally outperforms the descriptors-only variant at \SI{40}{\degreeCelsius} ($R^2 = 0.880$ vs.\ $0.872$), and at \SI{100}{\degreeCelsius} the descriptors-only model actually performs slightly better ($R^2 = 0.814$ vs.\ $0.807$). This suggests that fingerprints contribute little additional signal when descriptors are already present, and in some cases the additional dimensionality may introduce noise that slightly degrades performance. The SHAP analysis (Section~\ref{sec:shap_discussion}) provides further insight into which specific descriptors are most influential.

\subsection{MLP Failure on High-Dimensional Inputs}

The MLP (combined) and MLP (fingerprints) models fail catastrophically, despite the MLP (descriptors) variant performing well. This highlights a fundamental limitation of simple feed-forward neural networks on small datasets: the 2048-dimensional fingerprint input creates an extremely high-dimensional space that the MLP cannot navigate effectively with the available training data. The training curves confirm this interpretation --- the large train--validation loss gap for these models indicates severe overfitting, where the network memorises training examples rather than learning generalisable patterns.

XGBoost handles the same high-dimensional fingerprint input more gracefully ($R^2 = 0.526$) because its tree-based structure inherently performs feature selection (each split uses only one feature) and is regularised through subsample ratios and maximum tree depth. However, even XGBoost cannot extract strong predictive signal from fingerprints alone, reinforcing the conclusion that the fingerprint representation simply does not encode the physicochemical information most relevant to FOM1.

\subsection{ChemBERTa Embeddings: Strong Without Fine-Tuning}

A noteworthy result is that Ridge regression on frozen ChemBERTa embeddings ($R^2 = 0.858$ at \SI{40}{\degreeCelsius}, $0.823$ at \SI{100}{\degreeCelsius}) performs comparably to --- and at \SI{100}{\degreeCelsius} outperforms --- the fine-tuned ChemBERTa-2 model ($R^2 = 0.853$ and $0.766$ respectively). This suggests that ChemBERTa's pre-trained representations already encode much of the information needed for FOM1 prediction, and that fine-tuning on this relatively small dataset provides limited benefit, possibly even degrading performance through overfitting of the transformer backbone.

This has practical implications: extracting frozen embeddings and fitting a simple linear model takes under 2 seconds per fold, compared to over 800 seconds for fine-tuning, making it approximately 400 times faster with comparable or superior accuracy.

\subsection{Temperature Dependence of Model Performance}

All models perform worse at \SI{100}{\degreeCelsius} than at \SI{40}{\degreeCelsius}. This likely reflects the greater complexity of high-temperature fluid behaviour, where non-ideal effects (molecular decomposition pathways, non-linear viscosity--temperature relationships) become more pronounced and harder to predict from room-temperature molecular structure alone. The performance drop is largest for models that rely on flexible, data-hungry architectures (Deep Ensemble MLP drops by 0.148, ChemBERTa-2 by 0.087) and smallest for simpler models operating on informative features (Ridge embeddings drops by only 0.035, XGBoost descriptors by 0.058).

The change in ranking between temperatures --- particularly Ridge embeddings rising from fourth at \SI{40}{\degreeCelsius} to first at \SI{100}{\degreeCelsius} --- suggests that simpler models generalise more robustly when the prediction task becomes harder, likely because they are less prone to overfitting subtle patterns in the training data that do not transfer to the more complex high-temperature regime.

\subsection{BiLSTM Instability}

The BiLSTM exhibits high fold-to-fold variability (standard deviation of 0.334 and 0.280 in $R^2$ at \SI{40}{\degreeCelsius} and \SI{100}{\degreeCelsius} respectively), indicating that it is highly sensitive to the specific training data in each fold. This instability likely arises because character-level SMILES processing is inherently challenging: SMILES strings are not designed as a natural language and their syntax (e.g., ring closure numbers, branch parentheses) encodes structural information in ways that are difficult for sequential models to decode reliably, particularly from small datasets. The noisy training curves corroborate this interpretation.

\subsection{Computational Efficiency}

There is considerable variation in training time across architectures. Ridge regression on ChemBERTa embeddings completes in under 2 seconds per fold (including embedding extraction), while ChemBERTa-2 fine-tuning requires approximately 4500 seconds (75 minutes). XGBoost models fall in between, with training times of 150--330 seconds per fold depending on the feature set. Given that Ridge embeddings achieves competitive or superior accuracy to ChemBERTa-2 fine-tuning at a fraction of the computational cost, it represents an excellent accuracy--efficiency trade-off for this task.

\subsection{Interpretability Through SHAP}
\label{sec:shap_discussion}

The SHAP analysis reveals which molecular properties drive high cooling performance and, critically, \emph{in which direction}. The beeswarm plots (Figures~\ref{fig:shap_beeswarm_40} and~\ref{fig:shap_beeswarm_100}) make this directly readable: each point is one molecule, its horizontal position shows how much that feature pushes the FOM1 prediction up (right) or down (left), and its colour shows whether that molecule has a high (red) or low (blue) value of the feature. A cluster of red points on the left means that a high value of that feature \emph{decreases} FOM1; a cluster of red points on the right means it \emph{increases} FOM1.

\paragraph{What makes a molecule a good coolant?}
The beeswarm plots reveal a clear and physically intuitive picture of what a high-FOM1 molecule looks like:

\begin{enumerate}
    \item \textbf{Small molecules are better coolants.} The most important descriptor at both temperatures is \texttt{Chi0}, the Kier--Hall zeroth-order connectivity index, which is a direct proxy for molecular size (it sums over all bonds in the molecule). The beeswarm plots show an unambiguous pattern: molecules with \emph{high} \texttt{Chi0} (red points) have strongly \emph{negative} SHAP values --- they are pushed toward lower FOM1. Conversely, molecules with \emph{low} \texttt{Chi0} (blue points, i.e.\ small molecules) are pushed toward higher FOM1. The same pattern holds for \texttt{Chi0v} (which additionally weights by valence electrons), \texttt{Chi1n} (first-order connectivity), and \texttt{MolWt} (molecular weight). The physical explanation is straightforward: smaller molecules have lower viscosity and higher thermal conductivity per unit volume, both of which directly increase FOM1. This is the single strongest signal in the data --- at \SI{100}{\degreeCelsius}, \texttt{Chi0} alone has a mean $|\text{SHAP}|$ of 7.95, nearly double its importance at \SI{40}{\degreeCelsius} (4.72), because viscosity becomes an even more dominant factor in the FOM1 expression at elevated temperatures.

    \item \textbf{Structurally simple, compact molecules are preferred.} \texttt{FpDensityMorgan1} (the ratio of unique Morgan fingerprint fragments at radius~1 to the number of heavy atoms) consistently ranks in the top three at both temperatures, and high values push FOM1 \emph{upward} (red points on the right). A high fingerprint density means the molecule packs many distinct local chemical environments into a small number of atoms --- in other words, it is \emph{compact and functionally rich} rather than repetitive. This favours small, heteroatom-containing molecules (e.g.\ short-chain alcohols, amines, ketones) over long-chain hydrocarbons where the same \textrm{-CH\textsubscript{2}-} unit repeats many times.

    \item \textbf{Less lipophilic surface area improves cooling.} \texttt{SlogP\_VSA4}, which measures the van der Waals surface area contributed by atoms in an intermediate lipophilicity range, ranks in the top four at both temperatures. The beeswarm pattern is clear: high \texttt{SlogP\_VSA4} (red) $\rightarrow$ negative SHAP (worse FOM1). Molecules with large patches of moderately hydrophobic surface area are poor coolants. This is physically consistent: lipophilic surfaces promote weak van der Waals interactions between molecules, leading to lower heat capacity and thermal conductivity compared to molecules with polar surfaces that form hydrogen bonds or dipole--dipole interactions.

    \item \textbf{Some hydrophobic character is beneficial --- but distributed, not concentrated.} \texttt{BCUT2D\_LOGPHI} (the highest eigenvalue of the Burden connectivity matrix weighted by atomic lipophilicity) shows the opposite pattern to \texttt{SlogP\_VSA4}: high values push FOM1 \emph{upward}. This apparent contradiction resolves when considering what each descriptor measures. \texttt{BCUT2D\_LOGPHI} captures whether the molecule has \emph{any} hydrophobic region in its topology (a single non-polar moiety suffices to raise the eigenvalue), while \texttt{SlogP\_VSA4} measures the \emph{total area} of moderately lipophilic surface. The combination suggests that ideal coolants have a small non-polar region (e.g.\ a short alkyl chain, keeping the molecule liquid and reducing melting point) but not extensive hydrophobic surface area. In practical terms: a short-chain alcohol or ether is better than a long-chain hydrocarbon.

    \item \textbf{High molecular weight atoms in moderation help.} \texttt{BCUT2D\_MWHI} (highest eigenvalue of the Burden matrix weighted by atomic mass) shows high values pushing FOM1 upward. This likely captures the presence of one or two heavier heteroatoms (O, N, S) rather than overall molecular weight (which is penalised, as noted above). Heteroatoms increase polarity and hydrogen bonding capacity, improving heat capacity and thermal conductivity.
\end{enumerate}

\paragraph{Summary: the ideal coolant molecule.}
Combining these SHAP-derived insights, \textbf{a molecule with high FOM1 is small (low molecular weight, few bonds), compact rather than chain-like, with polar functional groups (e.g.\ hydroxyl, amine, carbonyl) that promote strong intermolecular interactions (hydrogen bonding, dipole--dipole), and only minimal hydrophobic surface area}. In practical chemical terms, this points toward molecules such as short-chain alcohols, polyols, amines, and small polar heterocycles as ideal cooling fluids. Large, branched, or predominantly non-polar molecules (long-chain alkanes, aromatic hydrocarbons) are predicted to be poor coolants. This is fully consistent with engineering experience: water and ethylene glycol --- small, highly polar molecules --- are among the best-known heat transfer fluids, while mineral oils and heavy hydrocarbons are far less effective per unit volume.

\paragraph{Temperature-dependent shifts.}
The cross-temperature comparison (Figure~\ref{fig:shap_cross_temp}) reveals that the same descriptors dominate at both temperatures, but their relative importance shifts. \texttt{Chi0} nearly doubles in importance from \SI{40}{\degreeCelsius} to \SI{100}{\degreeCelsius} (mean $|\text{SHAP}|$: $4.72 \rightarrow 7.95$), indicating that molecular size becomes an even stronger determinant of cooling performance at high temperature. This is physically expected: viscosity increases more steeply with molecular size at elevated temperatures, making the size penalty harsher. \texttt{MaxAbsEStateIndex}, which encodes the most extreme electrotopological state in the molecule (typically a highly polar functional group such as \textrm{-OH} or \textrm{-NH\textsubscript{2}}), rises from tenth to fifth in importance at \SI{100}{\degreeCelsius}, suggesting that strongly polar groups become more critical for maintaining good heat transfer at high temperatures. Conversely, \texttt{Chi4n} (a higher-order connectivity index capturing four-bond branching patterns) drops substantially at \SI{100}{\degreeCelsius}, indicating that fine details of molecular branching matter less when bulk size and polarity dominate.

The interpretability afforded by SHAP analysis is a practical advantage of XGBoost on descriptors over black-box alternatives such as ChemBERTa embeddings, where the latent features lack direct physical meaning. These insights can directly guide molecular screening: candidates should be pre-filtered for low molecular weight, high fingerprint density, and the presence of polar functional groups to enrich for high-FOM1 molecules.

% ============================================================
\section{Conclusion}
% ============================================================

This study compared ten machine learning architectures for predicting the first figure of merit for cooling (FOM1) of organic molecules at \SI{40}{\degreeCelsius} and \SI{100}{\degreeCelsius}. The key findings are:

\begin{enumerate}
    \item \textbf{XGBoost on descriptors (with or without fingerprints) is the best-performing architecture at \SI{40}{\degreeCelsius}}, achieving $R^2 = 0.880$ when using both descriptors and fingerprints, and $R^2 = 0.872$ with descriptors alone. The marginal contribution of fingerprints is negligible.

    \item \textbf{Ridge regression on frozen ChemBERTa embeddings is the best-performing architecture at \SI{100}{\degreeCelsius}} ($R^2 = 0.823$), demonstrating that pre-trained molecular language model representations encode substantial physicochemical information without any task-specific fine-tuning. It is also approximately 400 times faster than fine-tuning the full transformer.

    \item \textbf{Molecular descriptors are far more informative than Morgan fingerprints for FOM1 prediction.} This holds consistently across all model families (XGBoost, MLP) and both temperatures, reflecting the fact that FOM1 depends on bulk thermophysical properties rather than specific substructural motifs.

    \item \textbf{Simple models on good features outperform complex models on poor features.} Ridge regression on 384-dimensional ChemBERTa embeddings outperforms a fine-tuned 77M-parameter transformer, and XGBoost on 200 descriptors outperforms deep neural networks on 2048-dimensional fingerprints.

    \item \textbf{All models degrade at \SI{100}{\degreeCelsius}}, with simpler models showing more graceful degradation, suggesting that high-temperature FOM1 prediction is an inherently harder task.
\end{enumerate}

\subsection{Recommendations}

For deployment in a molecular screening or optimisation pipeline:
\begin{itemize}
    \item Use \textbf{XGBoost (descriptors)} as the primary model for FOM1 prediction. It achieves near-best accuracy at both temperatures, requires no GPU, trains in minutes, and uses interpretable physicochemical features.
    \item Use \textbf{Ridge regression on ChemBERTa embeddings} as a complementary model, particularly at \SI{100}{\degreeCelsius} where it is the top performer. Its sub-second inference time makes it suitable for screening large molecular libraries.
    \item If uncertainty quantification is required, use the \textbf{Deep Ensemble MLP} (which provides calibrated epistemic and aleatoric uncertainty estimates), but note its reduced accuracy relative to the top models.
\end{itemize}

\subsection{Future Work: Molecular Optimisation via Bayesian Optimisation and Reinforcement Learning}

The surrogate models developed in this study --- particularly the XGBoost (descriptors) and Deep Ensemble MLP models --- are not merely predictive tools; they form the foundation for automated molecular optimisation pipelines. Two complementary strategies are proposed: Bayesian optimisation (BO) for efficient exploration of molecular space, and reinforcement learning (RL) for generative molecular design.

\subsubsection{Bayesian Optimisation}

Bayesian optimisation is a sample-efficient strategy for optimising expensive black-box objective functions. In the molecular context, the objective is to find molecules that maximise FOM1 (or a multi-objective function incorporating FOM1 at both temperatures, synthetic accessibility, and safety constraints). The trained surrogate model serves as the objective function, replacing the need for costly thermophysical property calculations or experimental measurements at each evaluation.

The BO loop proceeds as follows:
\begin{enumerate}
    \item \textbf{Surrogate model:} The Deep Ensemble MLP is particularly well-suited as the BO surrogate because it provides both a mean prediction $\hat{y}(\mathbf{x})$ and a calibrated uncertainty estimate $\sigma(\mathbf{x})$ for each candidate molecule. The epistemic uncertainty (disagreement between ensemble members) naturally quantifies model confidence, which is essential for balancing exploration and exploitation.

    \item \textbf{Acquisition function:} An acquisition function such as Expected Improvement (EI) or Upper Confidence Bound (UCB) uses the surrogate's mean and uncertainty to select the next molecule to evaluate. Molecules in under-explored regions of chemical space (high $\sigma$) are favoured alongside molecules with high predicted FOM1 (high $\hat{y}$), preventing the optimiser from converging prematurely to a local optimum.

    \item \textbf{Molecular representation:} Candidate molecules are represented as descriptor vectors (the same 105 RDKit descriptors used in this study), and the BO search operates over this continuous descriptor space. A key challenge is the \emph{inverse mapping problem}: not every point in descriptor space corresponds to a valid molecule. This can be addressed by (i)~constraining the search to a pre-enumerated virtual library (e.g., GDB-17, ZINC), where BO selects the most promising candidates for evaluation; or (ii)~coupling BO with a generative model that maps from latent space to valid SMILES strings.

    \item \textbf{Iterative refinement:} After evaluating the selected molecule (either by computing its true FOM1 from thermophysical property predictions or by experimental measurement), the surrogate model is updated with the new data point, and the loop repeats. Typically 50--200 iterations suffice to identify high-performing candidates from libraries of millions.
\end{enumerate}

The XGBoost model, while lacking native uncertainty quantification, can also serve as a BO surrogate by bootstrapping predictions across the five cross-validation folds to estimate uncertainty, or by using a separate Gaussian process fitted to the XGBoost predictions.

\subsubsection{Reinforcement Learning for Generative Molecular Design}

While BO searches within a fixed molecular library, reinforcement learning (RL) can \emph{generate} entirely new molecular structures optimised for cooling performance. In this framework, a generative model (the \emph{agent}) learns to construct molecules token-by-token (as SMILES strings) or fragment-by-fragment, guided by a reward signal derived from the trained surrogate model.

The RL pipeline consists of the following components:
\begin{enumerate}
    \item \textbf{Policy network (agent):} A recurrent neural network (RNN) or transformer pre-trained on a large corpus of valid SMILES strings (e.g., ChEMBL, ZINC) to learn the grammar of molecular notation. This pre-training ensures that the agent generates syntactically valid and chemically plausible molecules from the outset.

    \item \textbf{Reward function:} The trained surrogate model (XGBoost or Deep Ensemble MLP) evaluates each generated molecule and returns a reward proportional to its predicted FOM1. The reward function can be augmented with additional terms to penalise undesirable properties:
    \begin{equation}
        R(\text{mol}) = w_1 \cdot \widehat{\text{FOM1}}_{40}(\text{mol}) + w_2 \cdot \widehat{\text{FOM1}}_{100}(\text{mol}) - w_3 \cdot \text{SA}(\text{mol}) - w_4 \cdot \mathbb{1}[\text{toxic}(\text{mol})]
    \end{equation}
    where SA is a synthetic accessibility score and the weights $w_i$ encode design priorities.

    \item \textbf{Policy gradient training:} The agent is fine-tuned using a policy gradient algorithm such as REINFORCE or Proximal Policy Optimisation (PPO). At each step, the agent generates a batch of molecules, computes their rewards, and updates its policy to increase the probability of generating high-reward structures. An entropy bonus or KL penalty relative to the pre-trained prior prevents mode collapse (the agent generating only slight variations of a single molecule).

    \item \textbf{Experience replay and diversity:} To maintain chemical diversity, techniques such as experience replay (storing and re-sampling high-reward molecules) and novelty bonuses (rewarding structural dissimilarity from previously generated molecules) can be employed.
\end{enumerate}

\subsubsection{Advantages of the Trained Surrogate Models}

The deep learning models developed in this study offer several specific advantages for BO and RL molecular optimisation:

\begin{itemize}
    \item \textbf{Fast inference:} Both XGBoost and the MLP ensemble evaluate a candidate molecule in under 1\,ms, enabling the reward function to be called millions of times during RL training or thousands of times per BO iteration without computational bottleneck. By contrast, computing FOM1 from first-principles thermophysical property models (group contribution methods, molecular dynamics) takes seconds to hours per molecule.

    \item \textbf{Differentiability:} The neural network models (MLP, Deep Ensemble, ChemBERTa) are fully differentiable, enabling gradient-based optimisation in continuous molecular representations. This allows the use of more sophisticated BO acquisition function optimisation and enables direct backpropagation of the reward signal through the surrogate in RL, which can accelerate convergence compared to score-function gradient estimators.

    \item \textbf{Uncertainty quantification:} The Deep Ensemble MLP provides calibrated uncertainty estimates that serve dual purposes: (i)~as the exploration term in BO acquisition functions, and (ii)~as a confidence filter in RL, where generated molecules with high surrogate uncertainty can be flagged for explicit validation rather than trusted blindly. This prevents the optimiser from exploiting regions where the surrogate is unreliable --- a well-known failure mode in model-based optimisation known as \emph{reward hacking}.

    \item \textbf{Multi-temperature optimisation:} Separate models at \SI{40}{\degreeCelsius} and \SI{100}{\degreeCelsius} enable Pareto-front optimisation, identifying molecules that perform well across the full operating temperature range rather than at a single condition. This is straightforward to implement in both BO (using multi-objective acquisition functions such as Expected Hypervolume Improvement) and RL (using a weighted multi-objective reward).
\end{itemize}

\subsubsection{Proposed Implementation}

We propose a two-stage optimisation pipeline:
\begin{enumerate}
    \item \textbf{Stage 1 --- Bayesian optimisation over virtual libraries.} Screen a large virtual library (e.g., the ZINC database of commercially available molecules, or an enumerated combinatorial library of candidate coolant scaffolds) using BO with the Deep Ensemble MLP as the surrogate. This identifies the highest-FOM1 molecules among known, synthesisable compounds and provides a performance baseline. The SHAP-derived insights (Section~\ref{sec:shap_discussion}) can be used to pre-filter the library, retaining only molecules with descriptor values in ranges associated with high FOM1 (e.g., low \texttt{Chi0}, moderate \texttt{SlogP\_VSA4}).

    \item \textbf{Stage 2 --- Reinforcement learning for de novo design.} Train a SMILES-based generative agent (pre-trained on ChEMBL) using the XGBoost surrogate as the reward function, with the BO-identified top molecules seeded into the replay buffer as starting points. The RL agent explores beyond the known chemical space, proposing novel molecular structures that may exceed the performance of any molecule in the virtual library. Generated candidates are validated by (i)~checking surrogate uncertainty (rejecting molecules where the Deep Ensemble disagrees strongly), (ii)~verifying synthetic accessibility, and (iii)~computing FOM1 from independent thermophysical property models for the top candidates.
\end{enumerate}

This two-stage approach combines the sample efficiency of BO (exploiting the known chemical space) with the creative capacity of RL (exploring beyond it), while the trained surrogate models provide the fast, accurate objective evaluations that make both approaches computationally tractable.

\end{document}
